% Claim statement file for row claim_3_24 -- "SigmaSeparationEquivalences".
% Auto-generated stub by scaffold/solve_chapter.py at row start. The
% manager agent (or a worker it dispatches) fills the body below with the
% claim's statement(s) from the lecture notes -- statement only, no proof
% (proofs live in the sibling `claim_3_24_proof_*.tex` file).
%
% This file is a sibling of `main.tex` inside `<subsection>/tex/`. The
% `subfiles` package lets it render both standalone and as part of
% `main.tex`. Note: `main.tex` does NOT \subfile this statement file -- the
% sibling `_proof_` file restates the statement above the proof, so
% including both in `main.tex` would be redundant. This file exists for
% standalone reading / grep convenience.

\documentclass[main]{subfiles}

\begin{document}
% ref: claim_3_24 (statement)
% title: SigmaSeparationEquivalences
\def\rowref{claim\_3\_24}
\phantomsection\label{claim_3_24}

\begin{Rem}[SigmaSeparationEquivalences]
    \begin{enumerate}
        \item By Proposition \ref{prp:sigma_opens} we have that $A \isPerp_G B \given C$ is equivalent to either of the following:
  \begin{enumerate}
      \item every walk from a node in $A$ to a node in $J \cup B$ is $C$-$\sigma$-blocked by $C$;
      \item every path from a node in $A$ to a node in $J \cup B$ is $C$-$\sigma$-blocked by $C$.
  \end{enumerate}
    \item Proposition \ref{prp:sigma_opens} also shows that if $A \nisPerp_G B \given C$ holds then:
  \begin{enumerate}
      \item there exists a (shortest) $C$-$\sigma$-open path from a node in $A$ to a node in $J \cup B$;
      \item there exists a (shortest) $C$-$\sigma$-open walk from a node in $A$ to a node in $J \cup B$ such that all its colliders lie in $C$.
  \end{enumerate}
  \end{enumerate}
  In practice we usually check if every path is $C$-$\sigma$-blocked or not. This is because there are, in contrast to walks, only a finite number of paths in a (finite) graph. In proofs, though, it often is easier to make use of walks, since these can be concatenated into walks (while one cannot in general concatenate two paths and again obtain a path).
\end{Rem}

\end{document}
